% Shared preamble for the English summaries in this repository.
\documentclass[11pt,a4paper]{article}
\usepackage[margin=2.2cm]{geometry}
\usepackage[T1]{fontenc}
\usepackage[utf8]{inputenc}
\usepackage{lmodern}
\usepackage{amsmath,amssymb}
\usepackage{graphicx}
\usepackage{booktabs}
\usepackage{enumitem}
\usepackage[hidelinks]{hyperref}
\setlist{itemsep=2pt,topsep=4pt}
\setlength{\parindent}{0pt}
\setlength{\parskip}{5pt}
\graphicspath{{figures/}}
\newcommand{\summaryheader}[3]{%
  \begin{center}
    {\Large\bfseries #2}\\[3pt]
    {\large #1}\\[3pt]
    {\small Arshia Goshtasbi \textbullet{} MSc Electrical Engineering (Control), Iran University of Science and Technology \textbullet{} #3}
  \end{center}
  \vspace{2pt}\hrule\vspace{8pt}
  \noindent\textit{English summary of my original Persian report, which is kept as a PDF in this folder.}\par}

\begin{document}
\summaryheader{Multivariable Control, Homework 4}{Poles of Multivariable Systems}{Autumn 1403 (2024)}

\begin{itemize}
  \item \textbf{The pole-polynomial rule.} I found the poles of
  \[
  G(s)=\frac{1}{(s+1)(s+3)^2}\begin{bmatrix}s^2+6s+9 & -1\\ 1 & s^2-9\\ 2(s^2-9) & 1\end{bmatrix}
  \]
  by reducing each entry, listing all non-zero minors, and taking the least common denominator. The pole polynomial is $(s+1)(s+3)^2$ (slide 17).
  \item \textbf{A second transfer matrix} (slide 27): its poles and their multiplicities.
  \item \textbf{Comparing pole-finding methods} (slide 47), from four representations:
    \begin{enumerate}
      \item state space: the eigenvalues of a minimal realisation;
      \item the Rosenbrock system matrix: watch for cancellations;
      \item the transfer matrix: the LCM of the denominators of all minors;
      \item matrix-fraction descriptions: the determinant of the denominator matrix.
    \end{enumerate}
\end{itemize}
\end{document}
